\section{Background and Motivation}
\label{sec:background}

\subsection{The Action-Only Execution Boundary}

An embodied \vla deployment has three logical roles.  An operator or scheduler
supplies a task, a policy maps its task view and multimodal observations to a
continuous \actionblock, and a consumer sends the block to an actuator-facing
controller.  After the block is consumed, the environment supplies a new
observation and the policy replans.  The policy therefore exports neither a
trusted task proof nor a predeclared reference trajectory; it exports an ordered
numeric proposal~\cite{pi05}.

For policy call $t$, write the source block as
\begin{equation}
A_t=(a_{t,0},\ldots,a_{t,H-1})\in\mathbb{R}^{H\times d}.
\label{eq:action-block-background}
\end{equation}
The consumer must decide whether this block is admissible, ensure that its rows
reach the controller in the intended order, and interpret the effects observed
after execution.  These are different questions: a task-relative assessment
does not itself establish execution integrity, and faithful dispatch does not
itself establish a physical safety property.

\subsection{How Attacks Reach the Same Object}

SABER perturbs policy-facing instructions, FreezeVLA perturbs visual input, and
embodied-agent jailbreaks manipulate language or multimodal context
~\cite{saber2026,freezevla2025,badrobot2025,robopair2025}.  Although these attacks
enter through different modalities, their physical influence is mediated by the
online block produced at the action-only interface.  They instantiate the
\intentaction surface; RQ1 exercises the SABER instruction channel.

After inference, a software
adversary can target the same lifecycle by substituting the source block,
reordering or omitting steps, replaying an earlier permission, or attaching
receipts and effects from a different proposal
~\cite{diat2019,cfaplus2024,arto2026,tat2026}.  These faults target the
\actionexecution transaction path enforced by \ltwoa.

The meaning of an accepted source action also depends on state and execution
configuration.  The same controller input can realize differently as joint
margins, contacts, or controller constraints change
~\cite{vlmpc2024,safetychance2024,measurementrobustcbf2021,
realizableshields2025}.  This state-dependent realization risk is the registered
target of \ltwob and need not arise from a separate attacker action.  The
security question spans the path from task authority, through the generated
block, to ordered execution and state-conditioned physical realization.

\subsection{Where Existing Protection Boundaries Stop}

Upstream defenses establish trusted inputs, protected digital control flow, or
an assessment over one or more policy candidates
~\cite{struq2025,secalign2025,camel2026,ace2026,sealvla2026,covervla2026}.
Their evidence ends before the full actuator transaction.  Execution-side
systems establish integrity for protected programs, control flow, or reference
motion~\cite{diat2019,cfaplus2024,arto2026,scaphy2023,tat2026}; their authority
begins from an already specified program or trajectory.  Prediction and
shielding qualify possible consequences or a local safety set, but do not by
themselves connect that decision to an independently anchored task assessment of
the newly generated block.

Table~\ref{tab:capability-comparison} summarizes these stopping points using
the six capabilities that define the protected lifecycle in this paper.

\begin{table*}[t]
\centering
\caption{Capability coverage across adjacent defenses at the action-only
boundary.  For a family row, a check means that at least one cited member states
a mechanism matching the column; a cross marks a capability outside the row's
stated endpoint.  The symbols do not rank overall system strength.}
\label{tab:capability-comparison}
\footnotesize
\setlength{\tabcolsep}{3.4pt}
\renewcommand{\arraystretch}{1.15}
\begin{tabularx}{\textwidth}{>{\raggedright\arraybackslash}X*{6}{>{\centering\arraybackslash}p{0.096\textwidth}}}
\toprule
Work family & \shortstack{Trusted/\\untrusted\\authority split}
  & \shortstack{Task--action\\assessment}
  & \shortstack{Exact action\\identity}
  & \shortstack{Receipt/effect\\binding}
  & \shortstack{Physical\\qualification}
  & \shortstack{Formal/\\attested\\semantics} \\
\midrule
StruQ, SecAlign~\cite{struq2025,secalign2025}
  & \pacovered & \panotcovered & \panotcovered & \panotcovered & \panotcovered & \panotcovered \\
AttriGuard~\cite{attriguard2026}
  & \pacovered & \pacovered & \pacovered & \panotcovered & \panotcovered & \panotcovered \\
CaMeL, ACE~\cite{camel2026,ace2026}
  & \pacovered & \pacovered & \pacovered & \panotcovered & \panotcovered & \pacovered \\
SEAL, CoVer~\cite{sealvla2026,covervla2026}
  & \panotcovered & \pacovered & \panotcovered & \panotcovered & \panotcovered & \panotcovered \\
DIAT, CFA+, ARTO, SCAPHY~\cite{diat2019,cfaplus2024,arto2026,scaphy2023}
  & \pacovered & \panotcovered & \panotcovered & \pacovered & \pacovered & \pacovered \\
TAT~\cite{tat2026}
  & \pacovered & \panotcovered & \panotcovered & \pacovered & \pacovered & \pacovered \\
Prediction and shielding~\cite{vlmpc2024,safetychance2024,measurementrobustcbf2021,realizableshields2025}
  & \panotcovered & \pacovered & \panotcovered & \panotcovered & \pacovered & \pacovered \\
\textbf{\system}
  & \pacovered & \pacovered & \pacovered & \pacovered & \pacovered & \pacovered \\
\bottomrule
\end{tabularx}
\end{table*}

The missing composition is therefore a proposal-scoped lifecycle with two
stages.  The first links authoritative task state to assessment of the source
block.  The second links an admitted block to ordered dispatch, receipts, and
effects, with an additional qualification obligation when robot state approaches
the registered boundary.  Section~\ref{sec:related} gives the detailed
comparison.  These stopping points yield the design requirements below.

\subsection{Three Design Challenges}

\textbf{C1. Bind semantic authority to a continuous proposal.}
A task phase can admit many valid continuous trajectories, so a monitor cannot
simply compare the policy output with one reference motion.  It needs an
independently anchored description of the current legal subtask and a checker
whose verdict states exactly which properties it covers.  Covered hard risks
must block execution, while uncertain task progress must not be presented as a
complete semantic guarantee.

\textbf{C2. Preserve proposal identity through execution.}
An assessment is useful only if later permission, dispatch, receipts, effects,
and task-phase advance refer to the same proposal and observation epoch.  The
consumer must reject substitution, replay, reordering, partial evidence, and
cross-proposal splicing while ensuring that failed evidence cannot be mistaken
for successful completion.

\textbf{C3. Qualify state-dependent realization without changing the source
action.}
Near a physical boundary, the controller configuration may need to change how a
source action is realized.  The design must keep the policy-produced row fixed,
make the selected guard/controller configuration explicit, and qualify that
combination before dispatch.  ``Same action'' therefore means the same policy
output at the protected controller boundary, not an assertion that every
controller configuration produces an identical physical trajectory.

\begin{table}[t]
\centering
\caption{The two alignment stages and their three component obligations.}
\label{tab:challenges}
\footnotesize
\begin{tabularx}{\columnwidth}{l>{\raggedright\arraybackslash}Xl}
\toprule
Component & Required invariant & Stage \\
\midrule
\lone (C1) & Verdict names the source block and checker scope & I \\
\ltwoa (C2) & Permission, dispatch, receipts, and effects share one proposal & II \\
\ltwob (C3) & Qualification names the fixed source action and configuration & II \\
\bottomrule
\end{tabularx}
\end{table}

The next section turns these challenges into an adversary model, trusted
assumptions, and lifecycle security objective.  Section~\ref{sec:design} then
presents the three corresponding components within the two stages.
